% Architecture Patterns Subsection
\subsection{Architecture and Algorithm Patterns}
\label{sec:architecture-patterns}

\subsubsection{Learning Paradigms}

Deep learning approaches in wearable seizure monitoring employ diverse learning paradigms. Feature-based approaches appear in four studies. These methods extract handcrafted features as input to classifiers. \textcite{fineDetectionKeyAutomated2025} extracted 594 features from 6-axis ACC and gyroscope data for an ANN classifier. \textcite{wangEpilepticSeizureDetection2025} engineered attitude angle features for an LSTM network. \textcite{singhrathoreDevelopmentMultimodalMachine2024} used multi-modal features for an MLP classifier.

End-to-end learning appears in two studies, both in forecasting. \textcite{meiselMachineLearningWristband2020} used an LSTM with 10 units operating directly on raw multi-modal wrist data. \textcite{nasseriAmbulatorySeizureForecasting2021} used a 4-layer LSTM with 128 hidden units receiving FFT channels of raw sensor data.

Hybrid approaches appear in two studies. \textcite{vielufSeizureMonitoringCombined2025} combined a DNN with handcrafted harmonic features extracted from 24-hour patterns. \textcite{elemamAutomatedValidatedTool2025} used separate CNNs for HRV and audio data without fusion between modalities.

Ensemble methods appear in detection applications. \textcite{spahrDeepLearningbasedDetection2025} trained 30 separate 1D CNN models with quantile aggregation. \textcite{dongTwoLayerEnsembleMethod2022} used a two-stage architecture with threshold-based pre-screening followed by a CNN-LSTM model.

Anomaly detection appears in two studies. \textcite{reintjesECGBasedDetectionEpileptic2025} compared three unsupervised anomaly detection methods on ECG data. \textcite{odeDevelopmentEpilepticSeizure2023} used a self-attentive autoencoder for patient-specific anomaly detection on ECG RRI intervals.

\subsubsection{Architecture Details}

Detailed descriptions of specific architectures are provided in Section~\ref{sec:modalities-architectures}. CNN-based approaches dominate detection research, with \textcite{spahrDeepLearningbasedDetection2025} implementing an ensemble of 30 CNN models with 14 convolutional layers each operating on ACC data sampled at 32~Hz with 30-second windows. LSTM architectures dominate forecasting research, with three of five forecasting studies using LSTM variants.

\subsubsection{Temporal Evolution}

Architecture choices have evolved over time. The earliest study, \textcite{borujenyDetectionEpilepticSeizure2013}, used traditional ANN approaches. Studies published between 2020 and 2022 focused on LSTM-based forecasting architectures. Studies published between 2023 and 2026 employed more complex architectures including ensembles, hybrid methods, and attention mechanisms.

\textbf{Key finding.} Detection research employs architectural diversity including CNNs, LSTMs, ANNs, and anomaly detection. Forecasting research concentrates heavily on LSTM variants. No forecasting study uses CNN-based approaches, possibly because forecasting requires modeling temporal dependencies rather than spatial patterns.

\textbf{Limitation.} Architecture descriptions vary in detail across studies. Some studies report complete architectural details while others provide minimal information. This variability complicates direct comparison and replication.
